\documentclass[10pt,landscape,a4paper]{article}
\usepackage[landscape,margin=1.2cm]{geometry}
\usepackage{multicol}
\usepackage{tabularx}
\usepackage{xcolor}
\usepackage{hyperref}
\usepackage[T1]{fontenc}
\usepackage{lmodern}
\usepackage{amsmath}
\usepackage{amssymb}
\usepackage{fuzz}
\usepackage{proof}

% Compact spacing
\setlength{\parindent}{0pt}
\setlength{\parskip}{2pt plus 1pt}
\setlength{\columnsep}{1.2cm}
\pagestyle{empty}

% Section formatting
\usepackage{titlesec}
\titleformat{\section}{\large\bfseries\color{darkgray}}{}{0pt}{}[\vspace{-4pt}\rule{\columnwidth}{0.4pt}]
\titlespacing{\section}{0pt}{6pt}{2pt}

% Code style
\newcommand{\kw}[1]{\texttt{\bfseries #1}}
\newcommand{\op}[1]{\texttt{#1}}
\newcommand{\renders}{\,$\rightarrow$\,}

% Stacked example: input then rendered output (compact)
\newenvironment{exwrite}%
  {\par\noindent\scriptsize\textbf{You write:}\par\vspace{0pt}\scriptsize}%
  {\vspace{1pt}}
\newenvironment{exget}%
  {\par\noindent\scriptsize\textbf{You get:}\par\vspace{0pt}\small}%
  {\par\vspace{3pt}}

% Table helper — left col monospace, right col roman
\newcolumntype{L}{>{\ttfamily}l}
\newcolumntype{R}{>{\raggedright\arraybackslash}X}

\begin{document}

\begin{center}
{\LARGE\bfseries txt2tex Reference}\\[2pt]
{\small Write math like you're at a whiteboard.  Get \LaTeX\ you can hand in.}\\[1pt]
{\footnotesize\color{gray}\url{https://github.com/jmf-pobox/txt2tex}}
\end{center}

\vspace{-2pt}

%% ── PAGE 1: FOUNDATIONS ──────────────────────────

\begin{multicols}{2}
\raggedcolumns

% ── DOCUMENT STRUCTURE ──────────────────────────
\section{Document Structure}

\textbf{Single-line directives.}

\begin{tabularx}{\columnwidth}{LR}
=== Title ===        & \verb|\section*{Title}| \\
** Solution N **     & Bold heading with spacing \\
(a) Text             & Part label \\
TITLE: My Doc        & Document title (\verb|\title{}|) \\
SUBTITLE: Sub        & Document subtitle \\
AUTHOR: Name         & Author name \\
DATE: 2026           & Date string \\
INSTITUTION: Univ    & Institution line \\
BIBLIOGRAPHY: f.bib  & BibTeX file \\
BIBLIOGRAPHY\_STYLE: plainnat & Citation style \\
CONTENTS:            & Table of contents \\
PARTS: inline        & Parts formatting style \\
PAGEBREAK:           & \verb|\newpage| \\
LINEBREAK:           & \verb|\medskip| (gap, no new page) \\
TEXT: prose here     & Smart prose (auto-detects formulas) \\
PURETEXT: raw \&\ text & Verbatim prose; escapes \LaTeX\ specials \\
LATEX: \textbackslash cmd & Raw \LaTeX\ passthrough (single-line) \\
{[cite key]}         & \verb|\citep{key}| (in TEXT blocks) \\
\end{tabularx}

\smallskip
\textbf{END-terminated blocks.}

\begin{tabularx}{\columnwidth}{LR}
LATEX:               & Multi-line raw \LaTeX\ block; body until \texttt{END} \\
B:                   & B-machine verbatim block; body until \texttt{END} \\
\end{tabularx}

{\footnotesize \kw{LATEX:} with content on the same line is single-line form.
\kw{LATEX:} alone on a line opens the block form — the body is live \LaTeX.
\kw{B:} wraps its body in \texttt{verbatim} — backslashes print as typed.
Both require a bare \texttt{END} at column~0 to close.}

\begin{exwrite}%
\begin{verbatim}
LATEX:
\begin{tikzpicture}
  \draw (0,0) -- (1,1);
\end{tikzpicture}
END
\end{verbatim}
\end{exwrite}

\begin{exwrite}%
\begin{verbatim}
B:
MACHINE Counter
VARIABLES n
INVARIANT n : NAT
END
\end{verbatim}
\end{exwrite}

\smallskip
\textbf{Identifier Decoration.}

\begin{tabularx}{\columnwidth}{LR}
count'               & $count'$ (after-state) \\
in?                  & $in?$ (input) \\
out!                 & $out!$ (output) \\
x?'                  & $x?'$ (combinations allowed) \\
'sunk'               & $\mbox{`sunk'}$ (string literal) \\
\end{tabularx}

{\footnotesize Decoration attaches to any identifier. Reserved words (\texttt{theta}, \texttt{defs}, \texttt{hide}, \texttt{project}, \texttt{Delta}, \texttt{Xi}) cannot be decorated.}

% ── PROPOSITIONAL LOGIC ─────────────────────────
\section{Propositional Logic}

\begin{tabularx}{\columnwidth}{LR}
p land q             & $p \land q$ \\
p lor q              & $p \lor q$ \\
lnot p               & $\lnot p$ \\
p => q               & $p \Rightarrow q$ \\
p <=> q              & $p \Leftrightarrow q$ \\
\end{tabularx}

{\footnotesize\textbf{Note:} Use \kw{land}, \kw{lor}, \kw{lnot}. English \texttt{and}/\texttt{or}/\texttt{not} are not supported.}

\smallskip
\textbf{Truth-table example:}

\begin{exwrite}%
\begin{verbatim}
TRUTH TABLE:
p | q | p => q
T | T | T
T | F | F
F | T | T
F | F | T
\end{verbatim}
\end{exwrite}
\begin{exget}%
\begin{center}
\begin{tabular}{ccc}
$p$ & $q$ & $p \Rightarrow q$ \\ \hline
T & T & T \\
T & F & F \\
F & T & T \\
F & F & T \\
\end{tabular}
\end{center}
\end{exget}

\textbf{Equivalence chain (\kw{EQUIV:}):}

\begin{exwrite}%
\begin{verbatim}
EQUIV:
lnot (p land q)
lnot p lor lnot q  [De Morgan]
\end{verbatim}
\end{exwrite}
\begin{exget}%
\begin{center}
$\begin{array}{l@{\hspace{2em}}r}
\lnot (p \land q) \\
\Leftrightarrow \lnot p \lor \lnot q
  & [\mbox{De Morgan}]
\end{array}$
\end{center}
\end{exget}

{\scriptsize \kw{EQUIV:} and its alias \kw{ARGUE:} join steps with $\Leftrightarrow$. Use for predicate equivalences.}

\textbf{Equality chain (\kw{EQUAL:}):}

\begin{exwrite}%
\begin{verbatim}
EQUAL:
0 + 0
0      [0+0 = 0]
length(<>)
       [length(<>) = 0]
\end{verbatim}
\end{exwrite}
\begin{exget}%
\begin{center}
$\begin{array}{l@{\hspace{1.5em}}r}
0 + 0 \\
= 0 & [\mbox{0+0 = 0}] \\
= length(\langle\rangle) & [\mbox{length}(\langle\rangle) = 0]
\end{array}$
\end{center}
\end{exget}

{\scriptsize \kw{EQUAL:} joins steps with $=$. Use when steps are arithmetic/set-valued, not propositional equivalences.}

% ── NUMBERS & ARITHMETIC ────────────────────────
\section{Numbers \& Arithmetic}

\begin{tabularx}{\columnwidth}{LR}
N                    & $\nat$ (natural numbers) \\
Z                    & $\num$ (integers) \\
N1                   & $\nat_1$ (positive naturals) \\
n + m / n - m        & addition / subtraction \\
n * m / n div m      & multiplication / integer division \\
n mod m              & modulo \\
n < m / n > m        & strict ordering \\
n <= m / n >= m      & non-strict ordering \\
n = m / n /= m       & equality / inequality \\
\end{tabularx}

{\footnotesize \textbf{Exponentiation:} \texttt{R\^{}n} (no space) emits \texttt{\textbackslash bsup n\textbackslash esup} — the fuzz \textit{iter} operator. It is only valid for relations, not arithmetic. For $x^2$ write \texttt{x * x}.}

\end{multicols}

%% ── PAGE 2: SETS, SEQUENCES, PREDICATE LOGIC ─────
\newpage

\begin{multicols}{2}
\raggedcolumns

% ── SETS & TYPES ────────────────────────────────
\section{Sets \& Types}

\begin{tabularx}{\columnwidth}{LR}
x elem S             & $x \in S$ \\
x notin S            & $x \notin S$ \\
A subset B           & $A \subseteq B$ \\
A psubset B          & $A \subset B$ (proper subset) \\
A union B            & $A \cup B$ \\
A intersect B        & $A \cap B$ \\
A setminus B         & $A \setminus B$ \\
bigcup S             & $\bigcup S$ \\
bigcap S             & $\bigcap S$ \\
power S              & $\power S$ \\
P1 S                 & $\power_1 S$ (non-empty) \\
F S                  & $\finset S$ \\
F1 S                 & $\finset_1 S$ (non-empty finite) \\
A cross B            & $A \cross B$ \\
\{x : S | P\}        & Set comprehension \\
n..m                 & $n \upto m$ (integer range) \\
\# S                 & $\# S$ (cardinality / length) \\
\end{tabularx}

{\footnotesize\textbf{Prefix-operator rule (\#133):} prefix-generic operators (\texttt{seq}, \texttt{P}, \texttt{dom}, \texttt{ran}, \texttt{bigcup}, …) require an atomic argument. The engine automatically wraps a nested prefix call in parentheses: \texttt{seq (P X)} $\to$ \verb|\seq~(\power X)|.}

\smallskip
\textbf{Set comprehension:}

\begin{exwrite}%
\begin{verbatim}
{x : N | x > 0 land x < 5}
\end{verbatim}
\end{exwrite}
\begin{exget}%
\begin{center}
$\{x : \nat \mid x > 0 \land x < 5\}$
\end{center}
\end{exget}

% ── PREDICATE LOGIC ─────────────────────────────
\section{Predicate Logic}

\begin{tabularx}{\columnwidth}{LR}
forall x : S | P     & $\forall x : S @ P$ \\
exists x : S | P     & $\exists x : S @ P$ \\
exists\_1 x : S | P  & $\exists_1 x : S @ P$ (unique) \\
lambda x : S | E     & $\lambda x : S @ E$ \\
mu x : S | P         & $\mu x : S @ P$ (definite desc.) \\
if p then e1 else e2 & conditional expression \\
\end{tabularx}

{\footnotesize \textbf{Bullet form:} \texttt{x : S | P} means $x : S \bullet P$ (bullet is implied). The vertical bar is the quantifier separator.}

\smallskip
\textbf{Tuple-pattern:}

\begin{exwrite}%
\begin{verbatim}
forall (x, y) : N cross N |
  x + y >= 0
\end{verbatim}
\end{exwrite}
\begin{exget}%
\begin{center}
$\forall (x, y) : \nat \cross \nat @ x + y \geq 0$
\end{center}
\end{exget}

{\footnotesize \textbf{Multi-decl:} separate variable groups with \texttt{;} inside a quantifier: \texttt{forall x : A; y : B | P}.}

{\footnotesize \textbf{Schema-text scope:} a bare \texttt{[d : T | P]} in a quantifier or function acts as an inline schema text.}

% ── RELATIONS ───────────────────────────────────
\section{Relations}

\begin{tabularx}{\columnwidth}{LR}
A <-> B              & $A \rel B$ (relation type) \\
x |-> y              & $(x \mapsto y)$ (maplet) \\
dom R / ran R        & $\dom R$ / $\ran R$ \\
id                   & identity relation \\
inv R / R\textasciitilde & $R^{-1}$ (inverse; postfix or prefix) \\
A dres R / A <| R    & $A \dres R$ (domain restrict) \\
A dsub R / A <<| R   & $A \ndres R$ (domain subtract) \\
R rres B / R |> B    & $R \rres B$ (range restrict) \\
R rsub B / R |>> B   & $R \nrres B$ (range subtract) \\
R oplus S / R ++ S   & $R \oplus S$ (relational override) \\
R o9 S               & $R \semi S$ (forward comp.) \\
R comp S             & $R \comp S$ (backward comp.) \\
R(| A |)             & $R\,[\,A\,]$ (relational image) \\
shows                & $\vdash$ (turnstile / judgment) \\
\end{tabularx}

{\footnotesize\textbf{o9 vs comp:} \texttt{o9} $\to$ \texttt{\textbackslash semi} (forward, $(R;S)(x)=S(R(x))$). \texttt{comp} $\to$ \texttt{\textbackslash comp} (backward, $(R\circ S)(x)=R(S(x))$).}

\smallskip
\textbf{Worked example — domain restrict:}

\begin{exwrite}%
\begin{verbatim}
{1, 2} dres {1 |-> a, 2 |-> b, 3 |-> c}
\end{verbatim}
\end{exwrite}
\begin{exget}%
\begin{center}
$\{1,2\} \dres \{1 \mapsto a, 2 \mapsto b, 3 \mapsto c\}$
\end{center}
\end{exget}

% ── FUNCTIONS ───────────────────────────────────
\section{Functions}

\begin{tabularx}{\columnwidth}{LR}
A +-> B              & $A \pfun B$ (partial) \\
A -> B               & $A \fun B$ (total) \\
A >+> B / A -|> B    & $A \pinj B$ (partial injection) \\
A >-> B              & $A \inj B$ (total injection) \\
A +->> B             & $A \psurj B$ (partial surjection) \\
A -->> B             & $A \surj B$ (total surjection) \\
A >->> B             & $A \bij B$ (bijection) \\
A 77-> B             & $A \ffun B$ (finite partial) \\
f(x)                 & function application \\
lambda x : T | f(x)  & $\lambda x : T @ f(x)$ \\
f ++ g               & $f \oplus g$ (override) \\
f~                   & $f^{-1}$ (postfix inverse) \\
\end{tabularx}

\smallskip
\textbf{Worked example — lambda:}

\begin{exwrite}%
\begin{verbatim}
double == lambda n : N | n * 2
\end{verbatim}
\end{exwrite}
\begin{exget}%
\begin{center}
$double == \lambda n : \nat @ n * 2$
\end{center}
\end{exget}

% ── SEQUENCES ───────────────────────────────────
\section{Sequences \& Bags}

\begin{tabularx}{\columnwidth}{LR}
<> / <a, b>          & $\langle\rangle$ / $\langle a, b \rangle$ \\
s \textasciicircum{} t  & $s \cat t$ (concat; space before \texttt{\^} required) \\
seq S / seq1 S       & $\seq S$ / $\seq_1 S$ \\
iseq S               & $\iseq S$ (injective sequences) \\
\# s                 & $\# s$ (length) \\
head s / tail s      & first / rest \\
last s / front s     & last / all-but-last \\
rev s                & reverse \\
s(i)                 & index (1-based) \\
s filter A           & $s \upharpoonright A$ (filter) \\
bag S                & $\bag S$ \\
{[}\![a, b]\!{]}     & $\lbag a, b \rbag$ (bag literal) \\
b1 bag\_union b2     & $b_1 \uplus b_2$ \\
b1 bag\_diff b2      & $b_1 \uminus b_2$ \\
\end{tabularx}

{\footnotesize\textbf{Concat:} write \texttt{<a> \^{} <b>} (space before \texttt{\^}).  No space $\to$ relation iteration (\texttt{\textbackslash bsup}). \textbf{bag\_diff:} emits \texttt{\textbackslash uminus} — bag subtraction clamped at zero.}

\end{multicols}

%% ── PAGE 3: Z PARAGRAPHS ─────────────────────────
\newpage

\begin{multicols}{2}
\raggedcolumns

% ── Z PARAGRAPHS ────────────────────────────────
\section{Z Paragraphs}

\begin{tabularx}{\columnwidth}{LR}
given A, B           & Given types: $[A, B]$ \\
T ::= a | b          & Free type definition \\
name == expr         & Abbreviation \\
\end{tabularx}

\smallskip
\begin{tabularx}{\columnwidth}{LR}
\multicolumn{2}{l}{\kw{axdef} / \kw{schema Name} / \kw{gendef [X]}}\\
\quad declarations   & Variable declarations \\
\texttt{where}       & Separator \\
\quad predicates     & Constraining predicates \\
\texttt{end}         & Close the block \\
\end{tabularx}

{\footnotesize \kw{axdef} introduces global constants; \kw{gendef [X]} parameterises by a type; \kw{schema Name} names a state or operation.}

\smallskip
\textbf{Zed block (unboxed paragraph):}

\begin{exwrite}%
\begin{verbatim}
zed
  given Entry
  Status ::= active | inactive
  MaxSize == 100
end
\end{verbatim}
\end{exwrite}
\begin{exget}%
\begin{zed}
[Entry] \also
Status ::= active | inactive \also
MaxSize == 100
\end{zed}
\end{exget}

{\scriptsize \kw{zed} blocks can mix given types, free types, abbreviations, and predicates in one environment. No visual box. Use for global type definitions; use \kw{axdef}/\kw{schema} for declarations with a \kw{where} clause.}

\smallskip
\textbf{Syntax block (multi-type free definitions):}

\begin{exwrite}%
\begin{verbatim}
syntax
  OP ::= plus | minus | times
  Expr ::= const<N> | binop<OP cross Expr>
end
\end{verbatim}
\end{exwrite}

{\scriptsize \kw{syntax} produces a column-aligned fuzz \texttt{syntax} environment. Blank lines between definitions emit \texttt{\textbackslash also} (visual group separator). Use for BNF-style grammars and type hierarchies.}

% ── SCHEMAS ─────────────────────────────────────
\section{Schemas}

{\scriptsize A named schema introduces a signature (declarations) and an invariant (where clause).}

\begin{exwrite}%
\begin{verbatim}
schema Counter
  size : N
  nums : seq N
where
  size = # nums
end
\end{verbatim}
\end{exwrite}
\begin{exget}%
\begin{schema}{Counter}
size : \nat \\
nums : \seq \nat
\where
size = \# nums
\end{schema}
\end{exget}

\smallskip
\textbf{Where-clause patterns:}

\textbf{1. Simple comparison.}
\begin{exwrite}%
\begin{verbatim}
schema Bounded
  size : N
where
  size < 10
end
\end{verbatim}
\end{exwrite}

\textbf{2. Set comprehension.}
\begin{exwrite}%
\begin{verbatim}
schema Inventory
  items : seq N
  stock : N
where
  stock = # {x : ran items | x > 0}
end
\end{verbatim}
\end{exwrite}

\textbf{3. Quantifier.}
\begin{exwrite}%
\begin{verbatim}
schema NonNeg
  s : seq Z
where
  forall i : 1..#s | s(i) >= 0
end
\end{verbatim}
\end{exwrite}

\smallskip
\textbf{Generic schemas.}

\begin{exwrite}%
\begin{verbatim}
gendef [X]
  empty : F X
where
  empty = {}
end
\end{verbatim}
\end{exwrite}
\begin{exget}%
{\scriptsize $\langle X \rangle$ is the type parameter; instantiate with \texttt{G[T]}.}
\end{exget}

\smallskip
\textbf{Schemas as predicates.}

{\scriptsize Inside a quantifier, a schema name acts as a predicate (both its signature and its where clause must be satisfied):}

\begin{exwrite}%
\begin{verbatim}
forall c : Counter |
  c.size >= 0
\end{verbatim}
\end{exwrite}
\begin{exget}%
\begin{center}
$\forall c : Counter @ c.size \geq 0$
\end{center}
\end{exget}

% ── SCHEMA OPERATIONS ───────────────────────────
\section{Schema Operations}

\textbf{Inclusion operators.}

\begin{tabularx}{\columnwidth}{LR}
SName                & bare inclusion \\
Delta Counter        & $\Delta Counter$ (before/after-state pair) \\
Xi Card              & $\Xi Card$ (read-only state) \\
theta S              & $\theta S$ (binding of current state) \\
theta S'             & $\theta S'$ (binding of after-state) \\
\end{tabularx}

\smallskip
\textbf{Horizontal definition.}

\begin{tabularx}{\columnwidth}{LR}
\multicolumn{2}{l}{\texttt{Name defs RHS}:}\\
Name defs Schema     & $Name \defs Schema$ \\
N defs Delta C       & $N \defs \Delta C$ \\
N defs [d : T | P]   & inline schema text \\
N[X] defs G[X]       & generic form \\
\end{tabularx}

\smallskip
\textbf{Schema calculus operators (RHS of \texttt{defs}).}

\begin{tabularx}{\columnwidth}{LR}
S[old/new]           & rename component \\
S[a/b, c/d]          & multiple renames \\
S ; T                & $S \semi T$ (composition) \\
S >> T               & $S \pipe T$ (piping) \\
S hide (x, y)        & $S \hide (x,y)$ \\
S project T          & $S \project T$ \\
\end{tabularx}

% ── Z BINDINGS ──────────────────────────────────
\section{Z Bindings}

{\scriptsize Z RM §3.7: binding brackets construct labelled tuples.}

\begin{tabularx}{\columnwidth}{LR}
\{| name == e |\}      & $\lblot name == e \rblot$ \\
\{| a == e, b == f |\} & multiple components \\
\{| |\}                & empty binding \\
\end{tabularx}

\smallskip
\textbf{Multi-typed comprehension with binding:}

\begin{exwrite}
\begin{verbatim}
{ t : Track; a : Album |
  t.albumId = a.albumId .
  {| trackId == t.trackId |} }
\end{verbatim}
\end{exwrite}
\begin{exget}
\begin{center}
$\{ t : Track; a : Album \mid t.albumId = a.albumId$\\
$\bullet \lblot trackId == t.trackId \rblot \}$
\end{center}
\end{exget}

{\scriptsize Use \texttt{;} to separate variable-type pairs inside \texttt{\{ \}}. Binding components are \textbf{comma}-separated (Z RM §3.7).}

\smallskip
{\scriptsize\textbf{Fuzz note:} bindings emit outside any Z environment. Do not place them inside \texttt{zed}/\texttt{axdef}/\texttt{schema} blocks — fuzz rejects \texttt{==} inside $\lblot \ldots \rblot$ in block contents.}

\end{multicols}

%% ── PAGE 4: PROOFS ───────────────────────────────
\newpage

\begin{multicols}{2}
\raggedcolumns

% ── PROOFS ──────────────────────────────────────
\section{Proofs}

\textbf{INFRULE: — named inference rule display:}

\begin{exwrite}%
\begin{verbatim}
INFRULE:
premise_1
premise_2
---
conclusion [Rule name]
\end{verbatim}
\end{exwrite}
\begin{exget}%
\begin{center}
$\infer[\mbox{Rule name}]{conclusion}{premise_1 & premise_2}$
\end{center}
\end{exget}

{\scriptsize \kw{INFRULE:} renders a labelled rule schema. Use for displaying rule definitions (not proofs). The three-dash line \texttt{---} separates premises from conclusion.}

\smallskip
\textbf{How proof trees work (\kw{PROOF:}):}

{\scriptsize The source is \textbf{conclusion-first} (rendered tree reads
bottom-up). Premises are indented under their conclusion. A rule with
two or more premises requires \texttt{::} on each premise --- without
it, indented lines collapse into a \emph{linear chain} (each line
becomes the sole premise of the line above) and the extra premises
are silently dropped. Unary rules take a single indented child with
no \texttt{::}.}

\smallskip
\textbf{Proof tree syntax:}

\begin{tabularx}{\columnwidth}{LR}
\textrm{indent}          & Premises indented under conclusion \\
{[rule]}                  & Justification label \\
{[rule from $n$]}         & Discharge assumption $n$ \\
{[$n$] expr [assumption]} & Boxed assumption labelled $n$ \\
expr [from $n$]           & Leaf reference to assumption $n$ \\
::                        & Premise of a multi-premise rule \\
case $X$:                 & Branch in a \kw{lor elim} case analysis \\
\end{tabularx}

{\footnotesize\textbf{Rules:} \texttt{land intro}, \texttt{land elim 1/2}, \texttt{lor intro 1/2}, \texttt{lor elim}, \texttt{=> intro/elim}, \texttt{<=> intro/elim}, \texttt{lnot intro/elim}, \texttt{forall intro/elim}, \texttt{exists intro/elim}.}

\smallskip
\textbf{Modus ponens (binary --- requires \texttt{::} on each premise):}

\begin{exwrite}%
\begin{verbatim}
PROOF:
q [=> elim]
  :: p => q
  :: p
\end{verbatim}
\end{exwrite}
\begin{exget}%
\begin{center}
$\infer[\Rightarrow\textrm{-elim}]{q}{p \Rightarrow q & p}$
\end{center}
\end{exget}

\textbf{$\land$-intro (binary):}

\begin{exwrite}%
\begin{verbatim}
PROOF:
p land q [land intro]
  :: p
  :: q
\end{verbatim}
\end{exwrite}
\begin{exget}%
\begin{center}
$\infer[\land\textrm{-intro}]{p \land q}{p & q}$
\end{center}
\end{exget}

\textbf{$\Rightarrow$-intro with discharge (use \kw{from N}):}

\begin{exwrite}%
\begin{verbatim}
PROOF:
p => p [=> intro from 1]
  [1] p [assumption]
  :: p [from 1]
\end{verbatim}
\end{exwrite}
\begin{exget}%
\begin{center}
$\infer[\Rightarrow\textrm{-intro}^{[1]}]{p \Rightarrow p}{\ulcorner p \urcorner^{[1]}}$
\end{center}
\end{exget}

{\scriptsize \kw{[1] p [assumption]} declares the discharge binder.
\kw{p [from 1]} marks every leaf use. \kw{[=> intro from 1]} discharges it.}

\textbf{$\lor$-elim / case analysis (ternary --- disjunction + two cases):}

\begin{exwrite}%
\begin{verbatim}
PROOF:
(p lor q) => r [=> intro from 1]
  [1] p lor q [assumption]
  :: r [lor elim from 1]
    :: p lor q [from 1]
    case p:
      :: r [...derive r from p...]
    case q:
      :: r [...derive r from q...]
\end{verbatim}
\end{exwrite}
\begin{exget}%
\begin{center}
$\infer[\lor\textrm{-elim}]{r}{p \lor q & \infer{r}{\ulcorner p \urcorner^{[1]}} & \infer{r}{\ulcorner q \urcorner^{[1]}}}$
\end{center}
\end{exget}

{\scriptsize \kw{lor elim from N} discharges the disjunction's assumption [N];
inside each \kw{case X:} block, the case formula is referenced as \kw{X [from N]}.}

\end{multicols}

%% ── PAGE 5: RELATIONAL DATABASE NOTATION ─────────
\newpage

\begin{multicols}{2}
\raggedcolumns

% ── RELATIONAL DATABASE NOTATION ────────────────
\section{Relational Database Notation}

\textbf{Primary key annotation.}

{\scriptsize Mark a primary-key attribute with \kw{pk} in a named schema body. The annotation sits outside the Z box so fuzz type-checks cleanly.}

\begin{exwrite}
\begin{verbatim}
schema Book
  pk bookId : BookId
  isbn      : ISBN
  pages     : N
end
\end{verbatim}
\end{exwrite}
\begin{exget}
\begin{center}
\begin{schema}{Book}
bookId : BookId \\
isbn : ISBN \\
pages : \nat
\end{schema}\\[2pt]
\noindent$\mathrm{PK}(\mathrm{Book}) = \{bookId\}$
\end{center}
\end{exget}

{\scriptsize Composite: two \kw{pk} lines $\Rightarrow$ $\mathrm{PK}(S)=\{a,b\}$.
Comma-separated names on one \kw{pk} line also work.
\kw{pk} is rejected in \texttt{axdef}/\texttt{gendef}/inline schema text.}

\smallskip
\textbf{Relational algebra.}

\begin{tabularx}{\columnwidth}{LR}
sigma[p](R)          & $\mathrm{Restrict}_{p}(R)$ — restrict \\
pi[A,B](R)           & $\mathrm{Project}_{A,B}(R)$ — project \\
rho[A as B](R)       & $\mathrm{Rename}_{A \to B}(R)$ — rename \\
R bowtie S           & $R \otimes S$ — natural join \\
R bowtie [p] S       & $\mathrm{Join}_{p}(R, S)$ — theta-join \\
R div S              & $R \div S$ — division \\
\end{tabularx}

{\footnotesize \texttt{sigma}/\texttt{pi}/\texttt{rho} bind at atom level. \texttt{bowtie} and \texttt{div} are infix at cross-product precedence.}

{\footnotesize\textbf{Fuzz note:} algebra expressions emit outside any Z environment — fuzz silently skips them. Do not wrap algebra inside \texttt{zed}/\texttt{axdef}/\texttt{schema} blocks.}

\smallskip
\textbf{FK predicate form.}

{\scriptsize Foreign-key constraints are expressed as axdef predicates:}

\begin{exwrite}
\begin{verbatim}
axdef
where
  (R, S) |-> {a |-> b} elem FK
end
\end{verbatim}
\end{exwrite}

\smallskip
\textbf{GROUP \& UNGROUP.}

{\scriptsize Date's nested-relation operators.}

\begin{tabularx}{\columnwidth}{LR}
R group (\{A,B\} as x) & $R\ \mathop{\mathrm{GROUP}} (\{A,B\}\ \mathop{\mathrm{AS}}\ x)$ \\
R ungroup x          & $R\ \mathop{\mathrm{UNGROUP}}\ x$ \\
\end{tabularx}

\begin{exwrite}
\begin{verbatim}
Members group ({name,email} as contact)
Members ungroup contact
\end{verbatim}
\end{exwrite}
\begin{exget}
\begin{center}
$Members\ \mathop{\mathrm{GROUP}}\ (\{name,email\}\ \mathop{\mathrm{AS}}\ contact)$\\[3pt]
$Members\ \mathop{\mathrm{UNGROUP}}\ contact$
\end{center}
\end{exget}

{\scriptsize\textbf{Fuzz note:} GROUP/UNGROUP emit outside any Z environment. Do not place them inside \texttt{zed}/\texttt{axdef}/\texttt{schema} blocks.}

\smallskip
\textbf{GROUP aggregator form.}

{\scriptsize The aggregator keywords are \kw{Count}, \kw{Sum}, \kw{Avg}, \kw{Min}, \kw{Max}, \kw{Median}. Each takes one attribute name and an alias. The regroup form (\texttt{\{...\} as alias}) and the aggregator form are mutually exclusive in a single \texttt{group} expression.}

\begin{tabularx}{\columnwidth}{LR}
R group (Count(x) as n)   & $R~\mathrm{Group}(\mathrm{Count}(x)~\mathrm{as}~n)$ \\
R group (Sum(x) as n)     & $R~\mathrm{Group}(\mathrm{Sum}(x)~\mathrm{as}~n)$ \\
R group (Avg(x) as n)     & $R~\mathrm{Group}(\mathrm{Avg}(x)~\mathrm{as}~n)$ \\
R group (Min(x) as n)     & $R~\mathrm{Group}(\mathrm{Min}(x)~\mathrm{as}~n)$ \\
R group (Max(x) as n)     & $R~\mathrm{Group}(\mathrm{Max}(x)~\mathrm{as}~n)$ \\
R group (Median(x) as n)  & $R~\mathrm{Group}(\mathrm{Median}(x)~\mathrm{as}~n)$ \\
\end{tabularx}

\begin{exwrite}
\begin{verbatim}
Sales group (Count(orderId) as orderCount,
             Sum(amount) as totalRevenue)
\end{verbatim}
\end{exwrite}
\begin{exget}
\begin{center}
$Sales~\mathrm{Group}(\mathrm{Count}(orderId)~\mathrm{as}~orderCount,$\\
$\mathrm{Sum}(amount)~\mathrm{as}~totalRevenue)$
\end{center}
\end{exget}

{\scriptsize Multiple aggregators are comma-separated. Each accepts exactly one attribute identifier — no nested aggregators, no compound expressions.}

\smallskip
\textbf{Quick reference.}

\begin{tabularx}{\columnwidth}{LR}
pk a : T             & primary key in schema \\
pk a, b : T          & composite pk \\
sigma[p](R)          & restrict \\
pi[A](R)             & project \\
rho[A as B](R)       & rename attr \\
R bowtie S           & natural join ($\otimes$) \\
R div S              & division \\
\{| a == e |\}       & binding bracket \\
\{S; C | P . E\}    & multi-typed set comp \\
R group (\{A\} as x) & regroup attrs \\
R group (Count(a) as n) & aggregate \\
R ungroup x          & ungroup attr \\
\end{tabularx}

\end{multicols}

\end{document}
